\documentclass[11pt]{article}

% Change "review" to "final" for the camera-ready version.
\usepackage[review]{acl}

\usepackage{times}
\usepackage{latexsym}
\usepackage[T1]{fontenc}
\usepackage[utf8]{inputenc}
\usepackage{microtype}
\usepackage{inconsolata}
\usepackage{graphicx}
\usepackage{booktabs}
\usepackage[table]{xcolor}
\usepackage{enumitem}
\usepackage{amsmath}
\usepackage{amssymb}
\usepackage{multirow}

\newcommand{\best}[1]{\textbf{#1}}
\newcommand{\code}[1]{\texttt{\small #1}}

\title{When the Label Is the Method:\\
An Audit of Scaffolding Gains in LLM Financial Extraction}

\author{Anonymous Submission}

\begin{document}
\maketitle

\begin{abstract}
% Why they should care, the problem.
Reading figures out of financial filings is a task where an error is invisible: a
misread number is a plausible number, and nothing in the output betrays it. Verifying
such extractions without a reference is therefore an attractive goal, and accounting
appears to supply the means, since financial statements obey arithmetic identities
that any correct reading must satisfy.
%
We built the system this invites: identity residuals as a reference-free detector,
a deterministic corrector, and a ladder of six configurations over two model scales,
evaluated on 104 annual reports against reference values taken from regulatory tags.
It worked: the corrector improved accuracy substantially over zero-shot extraction at
no additional inference cost, outperforming chain-of-thought self-verification at half
the model calls.
%
The improvement was not real. One extracted quantity is not covered by the reporting
standard, so its reference value has to be computed using the same formula the
corrector applies. Accuracy on that field therefore cannot fail, and it accounts for
the entire gain. Rescored without it, all six configurations are statistically
indistinguishable and four produce bit-identical values. The detector's apparent
precision has the same origin in a different form.
%
We give a general criterion for this failure, a one-line test that costs no
computation, and evidence that the inflated gain is predictable in advance from the
label definition alone. None of our scaffolding raises accuracy once the coupled field
is removed, and for the chain-of-thought configuration we show why: a verification
prompt can repair only the fields it is authorised to revise, so its coverage bounds its
gain regardless of its reasoning quality.
\end{abstract}

\section{Introduction}

Large language models are increasingly used where their output is a number rather than
a sentence. Financial due diligence, credit assessment and audit support all begin
with figures read out of filings \citep{wu2023bloomberggpt,10.1145/3768292.3770387},
and those figures propagate into every quantity computed from them.

This setting is unforgiving in a particular way. A hallucinated citation or an
incoherent summary can often be recognised from the text alone, and a body of work
detects exactly that \citep{ji2023survey,manakul2023selfcheckgpt}. A misread figure
offers no such signal: it is a well-formed number in a well-formed field, and it
differs from the correct one only in being wrong. Checking it appears to require the
reference value that production settings, by definition, do not have.

Accounting seems to offer an escape. Financial statements are not collections of
independent numbers but systems tied together by identities that hold by construction.
If an extraction violates one, at least one of the values involved is wrong, and this
is knowable without possessing any of them. The same identities can then repair what
they indict. It is a natural design, close in spirit to verifying generated code
against executable tests \citep{chen2023codet}, and this paper reports what happened
when it was built and then audited.

\paragraph{The system worked, and the result was an artifact.}
Evaluated against reference values drawn from the regulatory tags that accompany every
filing, the deterministic corrector improved field accuracy substantially over
zero-shot extraction, at no additional inference cost, and beat chain-of-thought
self-verification at half the model calls. Every one of those numbers is reproducible
from the released outputs, and all of them are artifacts of how one field was scored.

The mechanism is specific and it generalises beyond finance. EBITDA, one of the
extracted quantities, is not covered by the reporting standard, so no reference value
exists for it. The evaluation therefore constructs one, as operating income plus
depreciation and amortisation. The corrector's first rule computes that same field by
that same formula. The field is consequently scored correct precisely when two other
fields, each already scored on its own, were both correct. The evaluation counts one
success a second time under a third name, and the accounting identity it appears to
validate is the identity it was handed.

We state the condition generally.

\begin{quote}
\textbf{Method-coupled label.} A field whose reference value is produced by a formula
that the system under test also applies is method-coupled. Accuracy on such a field
cannot be falsified by the system's behaviour, and any gain attributed to it is
uninterpretable.
\end{quote}

The corresponding test is a single rescoring pass with coupled fields removed, which
requires no additional computation and no new data. Applied here it eliminates the
headline result entirely, on both model scales, with no residual. As
\S\ref{sec:predict} shows, the size of the inflation can be predicted in advance from
the label definition alone, before any rescoring is done.

Evaluation artifacts have repeatedly inflated apparent progress in NLP, through
annotation shortcuts \citep{gururangan2018annotation,poliak2018hypothesis}, benchmarks
solvable without the input \citep{kaushik2018much}, and test-set contamination
\citep{sainz2023nlp,recht2019imagenet}. The mechanism here is distinct and, we argue,
harder to notice: there is no leak and no shortcut, only a reference value that a rule of
the system reconstructs exactly, which survives a held-out split and any contamination
check. We do not claim that published systems share the defect, having audited only our
own; we claim that the condition is cheap to check and that we found no reason it should
be rare wherever a missing label is filled by computation.

The venue's open research policy encourages releasing code, prompts, configurations and
outputs so that findings can be re-evaluated later. This paper is that policy applied
to its own authors: every result below was obtained by rescoring already-released
outputs, offline, with no further inference.

\paragraph{Contributions.}
\begin{enumerate}[leftmargin=*, itemsep=2pt, topsep=2pt]
\item \textbf{A criterion, a test, and a predictor for method-coupled labels}
      (\S\ref{sec:audit}--\S\ref{sec:inject}), demonstrated on a case where coupling
      accounts for the whole of a reported improvement, where its magnitude is
      recoverable from the label definition alone, and where the mechanism reproduces
      under controlled injection across fifteen synthetic couplings.
\item \textbf{An exact null for these verification implementations, with a mechanism}
      (\S\ref{sec:collapse}): across six configurations and two model scales, none of
      our scaffolding raises accuracy, four configurations producing values identical
      field by field. The null is not evidence against the underlying techniques. For
      the chain-of-thought configuration we show why it holds here: a verification
      prompt can only repair what it is authorised to touch, and ours was scoped to a
      single identity. Verification coverage, rather than verification quality, is what
      bounds the achievable gain.
\item \textbf{An annotation-free evaluation protocol and a catalogue of the traps it
      exposes} (\S\ref{sec:protocol}, \S\ref{sec:traps}), covering label coupling,
      ablation by accuracy delta, threshold mismatch and unit scale, each stated as a
      mechanism with the instrument that detects it, and released with the artifacts
      needed to reproduce every table.
\end{enumerate}

\section{Related Work}

\paragraph{Financial document understanding.}
Numerical reasoning over financial text is served by FinQA \citep{chen2021finqa},
ConvFinQA \citep{chen2022convfinqa}, TAT-QA \citep{zhu2021tatqa}, FinanceBench
\citep{islam2023financebench} and DocMath-Eval \citep{zhou2023docmath}, aggregated by
FinBen \citep{xie2024finben}; BloombergGPT \citep{wu2023bloomberggpt} adapts a model to
the domain. These evaluate whether a model can \emph{compute} an answer from curated
evidence. Our question is upstream: whether a value read from a raw filing is the value
the filer reported. FiNER \citep{loukas2022finer} uses the same regulatory tags as
supervision for entity tagging; we use them as evaluation labels, and show that this
requires care whenever a label must be derived rather than read. Layout-aware
pretraining \citep{xu2020layoutlm} and table encoders \citep{herzig2020tapas} improve
the reading itself and are complementary to verification.

\paragraph{Reference-free verification.}
Detectors built on sampling disagreement \citep{manakul2023selfcheckgpt} or model
confidence \citep{kadavath2022know} estimate whether a model is uncertain, not whether
a value is right; a model can be perfectly consistent about the wrong table row.
Constraint-based verification is closer to program synthesis, where generated tests
\citep{chen2023codet} and grammar-constrained decoding
\citep{poesia2022synchromesh,willard2023efficient} check outputs against a
specification, and to weak supervision \citep{ratner2017snorkel}, where labelling
functions encode domain knowledge. Accounting identities are a strong instance of that
family. The failure we document applies to it generally: a labelling function that also
generates the label cannot test anything.

\paragraph{Self-correction.}
\citet{huang2024selfcorrect} and \citet{stechly2024selfverification} argue that
intrinsic self-correction fails without external feedback, and \citet{kamoi2024when}
survey when it succeeds. Chain-of-Verification \citep{dhuliawala2023chain} supplies
feedback through generated verification questions; our chain-of-thought configuration
is the same design applied to a single arithmetic identity. It appeared to help;
measured outside the coupled field it changes nothing.

\paragraph{Evaluation artifacts and reporting practice.}
Apparent progress has repeatedly been traced to properties of benchmarks rather than of
systems \citep{gururangan2018annotation,poliak2018hypothesis,kaushik2018much,
bowman2021will}, and contamination has become a standard confound to rule out
\citep{sainz2023nlp}. Behavioural testing \citep{ribeiro2020beyond}, careful
experimental reporting \citep{dodge2019show} and warnings about
explanation-after-the-fact \citep{lipton2019troubling} all respond to the same
pressure. Label coupling belongs on that list: it is invisible to contamination checks,
survives a held-out split, and is introduced not by the data but by a modelling decision
about how to fill a gap in the reference. Its systems analogue is the feedback loop
described by \citet{sculley2015hidden}, in which a component's output silently becomes
part of its own input.

\section{Protocol}
\label{sec:protocol}

\paragraph{Task.} Given the text of an annual report, extract five income-statement
fields (revenue, operating income, EBITDA, depreciation and amortisation, net income)
and three balance-sheet fields (total assets, total liabilities and equity, total
equity) for the most recent fiscal period. A field is correct when it falls within a
relative tolerance of the reference value; $\pm1\%$ is the operating point, and
\S\ref{sec:survives} reports the sweep.

\paragraph{Reference values.} Filings carry machine-readable tags, so the value behind
each reported figure is public. Fields are mapped onto standard concepts
(Appendix~\ref{app:concepts}) and the concept behind every value is recorded, making
each label auditable. Two properties of that mapping matter. Tag coverage is a property
of the filer, not of the corpus: total liabilities is tagged by roughly two thirds of
companies and operating income by four fifths, so the effective sample differs by field
and is reported per field. And EBITDA has no tag at all, being outside the
standard, so its reference value must be constructed. That decision is the subject of
\S\ref{sec:audit}.

\paragraph{Corpus.} From 62 large-cap tickers across 11 sectors we take the two most
recent annual reports, convert the primary document to text preserving table row
structure, and slice the section around the consolidated income statement. That heading
recurs in a filing, so we select the occurrence followed by the highest density of
printed money figures, which distinguishes the table from a reference to it without
per-filer rules. Each document's fiscal period is read from the filing's own report
date, so a document and its reference values cannot refer to different years. The
arithmetic closes exactly: $124$ possible documents, less $8$ filings that do not exist
in the recent window, less $12$ whose reporting scale cannot be read from their own
text, giving $104$ documents from $55$ companies (Appendix~\ref{app:corpus}).

\paragraph{Reporting scale.} Reference values are absolute amounts; statement tables
print thousands or millions. The multiplier is inferred from the filing text alone,
with the matched evidence recorded. The direction matters: snapping an extraction to
whichever power of a thousand lands nearest the reference would leak the label into the
prediction, which is the defect this paper is about.

The inference is itself validated, since everything downstream depends on it. For each
filing the ratio between a reference value and the value printed in its table is
observable, and a correct multiplier is the one that makes that ratio unity. Read this
way, the inference agrees with the observed multiplier on $101$ of $104$ filings
($97.1\%$); the three disagreements are documents where the model misread a value
rather than filings whose header was misread.

\paragraph{Configurations.} Six, sharing one extraction prompt so that differences are
attributable to the scaffolding: B1 a regex parser given the same patterns the
full system uses internally; B2 zero-shot extraction (one call); \textbf{B3}
B2 with self-consistency, five samples at $T{=}0.7$, field-wise median; B4 B2
with chain-of-thought self-verification of the EBITDA identity (two calls); \textbf{B5}
B2 with the detector, values untouched; B6 B5 with the corrector. The
corrector applies eight named rules to a fixed point, so its outcome is
order-independent and idempotent, and issues no model calls, since repairing an
extraction never requires re-extracting it. Models are an 8B and a 49B
instruction-tuned model on one hosted API, temperature $0$ except B3, seed fixed. Rate
limiting prevented completing the larger model on the full corpus, so scales are
compared on documents both completed; sample sizes accompany every table.

\begin{figure*}[t]
\centering
\includegraphics[width=\linewidth]{figures/tolerance_coupling.pdf}
\caption{\textbf{(a)} Accuracy depends strongly on the tolerance at which a field is
called correct, which is a free parameter of any such evaluation. \textbf{(b)} The
ladder as published (hatched) and rescored with the coupled field removed (solid).
The apparent progression from B2 to B6 is produced entirely by that one field; without
it every configuration lands on the same value, marked by the dotted line.}
\label{fig:ladder}
\end{figure*}

\section{The Result as It First Appeared}

Table~\ref{tab:asfirst} reports field accuracy over the five income-statement fields.
Identity repair gains $10.7$ and $13.0$ points over zero-shot extraction while issuing
no additional model calls; chain-of-thought verification gains $6.3$ and $12.7$ at
twice the calls. B5 equals B2 exactly, confirming that the detector judges without
altering values. Read on its own, the table says that cheap symbolic structure
substitutes for inference-time compute, and that it does so at both model scales.

\begin{table}[t]
\centering\small
\setlength{\tabcolsep}{4pt}
\begin{tabular}{@{}llrrrr@{}}
\toprule
& \textbf{Config.} & \textbf{Acc.} & \textbf{$\Delta$B2} & \textbf{$n$} & \textbf{Calls} \\
\midrule
\multirow{6}{*}{\rotatebox{90}{\textsc{8b}}}
& B1 regex            & 69.9 & $-3.6$ & 355 & 0 \\
& B2 zero-shot        & 73.5 & ref.   & 355 & 1 \\
& B3 self-cons.       & 72.7 & $-0.8$ & 355 & 5 \\
& B4 CoT verify       & 79.8 & $+6.3$ & 352 & 2 \\
& B5 ${+}$ detector   & 73.5 & $+0.0$ & 355 & 1 \\
& B6 ${+}$ corrector  & \best{84.2} & \best{$+10.7$} & 355 & 1 \\
\midrule
\multirow{5}{*}{\rotatebox{90}{\textsc{49b}}}
& B1 regex            & 69.9 & $-7.0$  & 355 & 0 \\
& B2 zero-shot        & 76.9 & ref.    & 355 & 1 \\
& B4 CoT verify       & 89.6 & $+12.7$ & 299 & 2 \\
& B5 ${+}$ detector   & 76.9 & $+0.0$  & 355 & 1 \\
& B6 ${+}$ corrector  & \best{89.9} & \best{$+13.0$} & 355 & 1 \\
\bottomrule
\end{tabular}
\caption{Field accuracy (\%) at $\pm1\%$, five income-statement fields, over $83$
documents. $n$ is the number of scored instances, which is lower for B4 where a
verification call did not return. Every figure is reproducible from the released
outputs.}
\label{tab:asfirst}
\end{table}

\section{The Audit}
\label{sec:audit}

\subsection{One field produces the entire gain}

EBITDA has no tag, so its reference value is constructed as operating income plus
D\&A, and the corrector's first rule computes EBITDA as operating income plus D\&A.
The field is therefore scored correct exactly when the extracted operating income and
D\&A are both correct, and both of those are already scored separately. Nothing about
the system's behaviour can make that field wrong while leaving its constituents right.

Rescoring with the coupled field excluded removes the effect entirely, on both scales,
with no residual (Table~\ref{tab:rescore}, Figure~\ref{fig:ladder}b). Every field that
the corrector turns from wrong to right is EBITDA: $38$ of $38$ on the smaller model,
$46$ of $46$ on the larger. Coupled instances are $15.5\%$ of those scored and $100\%$
of the gain.

\begin{table}[t]
\centering\small
\begin{tabular}{@{}lrrl@{}}
\toprule
\textbf{Model} & \textbf{with} & \textbf{without} & \textbf{fields repaired} \\
\midrule
8B  & $+10.7$ & \best{$+0.0$} & 38/38 EBITDA \\
49B & $+13.0$ & \best{$+0.0$} & 46/46 EBITDA \\
\bottomrule
\end{tabular}
\caption{Corrector gain over zero-shot extraction, with and without the coupled field.
The last column counts which fields the corrector repairs: on both scales, only the
coupled one.}
\label{tab:rescore}
\end{table}

A second defect compounds the first and is independent of it. The prompt instructs the
model to return a null value when a quantity is absent; EBITDA is absent from every
income statement in the corpus; the model complied in every case; and the metric scored
that as an error. Removing the field raises zero-shot accuracy by $13.5$ and $14.1$
points. The baseline was being penalised for following its instructions, which is most
of the room the corrector then appeared to recover.

\subsection{The inflation is predictable in advance}
\label{sec:predict}

Excluding a coupled field diagnoses the problem after the fact. The stronger question
is whether the inflation could have been anticipated. For a derived label
$L = g(f_1,\dots,f_k)$ and a system rule computing the same field by the same $g$, the
field is repaired exactly on documents where it was absent from the raw extraction and
every constituent was already correct. That count is computable from the raw extraction
and the rule definition, with no rescoring:
\[
\widehat{n} \;=\; \bigl|\{d : L \text{ absent in } d,\; f_i \text{ correct } \forall i\}\bigl|.
\]
On the smaller model this predicts $38$ repaired instances against $38$ observed, an
error of zero; on the larger, $43$ against $46$, an error of $6.5\%$, the discrepancy
arising from tolerance interactions when constituents are correct but near the
boundary.
A coupled field can therefore be flagged, and its contribution quantified, before any
experiment is rerun.

\subsection{The mechanism reproduces under controlled injection}
\label{sec:inject}

The EBITDA case is observational: that field happened to be untagged, a label had to be
constructed, and a rule happened to reconstruct it. To show the effect is a property of
the construction rather than of the field, we induce it deliberately.

For every pair $(a,b)$ of independently tagged fields we define a synthetic field $S$
with reference value $\mathrm{truth}(a) + \mathrm{truth}(b)$, declare it absent from
the raw extraction as an unreported quantity would be, and add the corrector rule
$S := \mathrm{extracted}(a) + \mathrm{extracted}(b)$. This is the EBITDA construction
with different operands. The phantom gain is the share of documents on which the
injected field is then scored correct.

Across fifteen injected couplings per model, phantom gain tracks the joint accuracy of
the constituents closely (Figure~\ref{fig:inject}; Pearson $r = 0.92$ and $0.89$, mean
absolute error $0.06$ and $0.05$), and ranges from $0.68$ to $1.00$. Two consequences
follow. First, a coupled label yields a free gain whose size is set by how good the
system already is on the parts, so the practice is most rewarding exactly where a
system is strongest, and therefore most tempting on a mature pipeline. Second, the
points lie slightly above the diagonal, because errors in the two constituents can
cancel in the sum: a coupled field is scored correct somewhat more often than all of
its constituents are, which makes the inflation marginally worse than the naive
prediction.

\begin{figure}[t]
\centering
\includegraphics[width=0.92\linewidth]{figures/coupling_injection.pdf}
\caption{Deliberately injected couplings. Each point is one synthetic field defined as
the sum of two independently tagged fields, with a corrector rule computing it the same
way. The phantom gain it produces is determined by the joint accuracy of its
constituents, irrespective of which fields are chosen.}
\label{fig:inject}
\end{figure}

\subsection{The ladder collapses}
\label{sec:collapse}

Excluding the coupled field leaves the configurations indistinguishable
(Table~\ref{tab:mcnemar}). Four of them (B2, B4, B5 and B6) produce values identical
field by field, not merely equal in aggregate: on the larger model a single instance out
of $300$ differs, and it does not change correctness. No adjacent pair shows a
significant difference, and the only significant comparison in the entire study is
between the language model and the pattern matcher.

\begin{table}[t]
\centering\small
\setlength{\tabcolsep}{4pt}
\begin{tabular}{@{}llrrr@{}}
\toprule
& \textbf{Pair} & \textbf{A only} & \textbf{B only} & \textbf{$p$} \\
\midrule
\multirow{5}{*}{\rotatebox{90}{\textsc{8b}}}
& B2 vs B3 &  9 &  6 & $0.607$ \\
& B2 vs B4 &  0 &  0 & $1.000$ \\
& B2 vs B5 &  0 &  0 & $1.000$ \\
& B2 vs B6 &  0 &  0 & $1.000$ \\
& B1 vs B2 & 14 & \best{60} & \best{$6{\cdot}10^{-8}$} \\
\midrule
\multirow{4}{*}{\rotatebox{90}{\textsc{49b}}}
& B2 vs B4 &  0 &  0 & $1.000$ \\
& B2 vs B5 &  0 &  0 & $1.000$ \\
& B2 vs B6 &  0 &  0 & $1.000$ \\
& B1 vs B2 &  8 & \best{66} & \best{$2{\cdot}10^{-12}$} \\
\bottomrule
\end{tabular}
\caption{Discordant pairs and exact McNemar tests, coupled field excluded, $300$
instances over $83$ documents. B4 is compared on the $297$ and $253$ instances it
returned.}
\label{tab:mcnemar}
\end{table}

Self-consistency is the one configuration that is not inert but merely ineffective. It
moves $38$ of $300$ values and changes $15$ correctness outcomes, nine in one direction
and six in the other ($p = 0.607$). The displacements are small and systematic: one filer's
revenue read as $57{,}400$ rather than $57{,}300$, another's D\&A as $52{,}785$ rather
than $52{,}795$, consistent with sampling changing which table row is read rather than
which value is inferred. Where the answer is a lookup, disagreement across samples
reflects retrieval noise, not reasoning uncertainty, and a median over five readings of
five different rows is not a better reading. This is a mechanism for the null on this
task, not a contradiction of self-consistency on reasoning benchmarks
\citep{wang2023selfconsistency}.

Chain-of-thought verification is inert for a reason internal to how it was specified,
and the reason is worth separating from the result. Its prompt authorises changes to
the coupled field only, so outside that field it cannot alter anything, and its apparent
property of never degrading a value is vacuous. This is not evidence that
chain-of-thought verification fails on this task; it is evidence that a verification
prompt bounds its own achievable gain by the set of fields it is allowed to revise. A
verifier scoped to one identity can repair at most the fields that identity relates,
whatever its reasoning quality. The practical reading is that verification coverage
should be reported alongside verification accuracy, since the first caps the second.

\subsection{Ablation deltas misdiagnose the rule set}

Leave-one-out ablation reports all eight correction rules as contributing nothing, on
both scales. We originally read that as mutual redundancy under fixed-point iteration.
Execution counts show otherwise (Table~\ref{tab:fires}): one rule fires, on every
document; six never fire; the eighth executes only when the first is removed. The regex
backfill fires zero and two times and never changes a scored outcome.

\begin{table}[t]
\centering\small
\begin{tabular}{@{}lrr@{}}
\toprule
\textbf{Rule} & \textbf{8B} & \textbf{49B} \\
\midrule
\code{ebitda\_from\_parts}       & \best{83} & \best{69} \\
\code{ebitda\_below\_ebit}       & 0 $\;$($78^{\dagger}$) & 0 $\;$($65^{\dagger}$) \\
other six rules, each            & 0 & 0 \\
\code{regex\_backfill}           & 0 & 2 \\
\bottomrule
\end{tabular}
\caption{Field instances changed, $83$ documents. $^{\dagger}$Executions when the first
rule is disabled. An accuracy delta of zero is consistent with redundancy and with
irrelevance; execution counts separate them.}
\label{tab:fires}
\end{table}

The point transfers to any rule set applied to convergence: report how often each rule
executes, because the two readings of a zero delta call for opposite responses.

\subsection{The detector is coupled in a different way}

\begin{figure}[t]
\centering
\includegraphics[width=\linewidth]{figures/threshold_surface.pdf}
\caption{Detector $F_1$ as a surface over the two thresholds an evaluation chooses
independently: the residual at which the detector fires ($\tau_d$) and the tolerance at
which a field counts as correct ($\tau_c$). Performance forms a ridge along
$\tau_d = \tau_c$. Choosing them independently, which is natural since they
belong to different parts of the pipeline, lands off the ridge and manufactures
errors that belong to the mismatch rather than to the detector.}
\label{fig:surface}
\end{figure}

The balance-sheet detector appeared to work. Firing on any nonzero residual of
$|\text{assets} - \text{liabilities and equity}|$, it reached a precision of $80.0\%$
(95\% CI $[56.3, 94.3]$) and a recall of $84.2\%$ over $68$ documents. Its four false
positives are all documents where both fields fall inside the $\pm1\%$ correctness
tolerance while the residual is nonzero. That is not a detector error but a threshold
mismatch, and Figure~\ref{fig:surface} shows its geometry: performance is maximal along
the diagonal where the two thresholds agree, and degrades away from it. At matched
thresholds precision is $100\%$ (CI $[79.4, 100]$) with no false positives.

That precision is nonetheless not a measurement. The two inputs of the identity map to
reference concepts that the standard defines to be equal, so a nonzero residual proves
an extraction error by construction. The detector was verifying that two printed lines
agree, which is useful but is not the identity claimed. Figure~\ref{fig:ecdf} shows the further
reason an area under a curve is the wrong summary here: most documents sit exactly at
zero residual, so any ranking metric is dominated by ties.

\begin{figure}[t]
\centering
\includegraphics[width=\linewidth]{figures/residual_ecdf.pdf}
\caption{Distribution of the identity residual. The mass at zero is what makes ranking
metrics misleading: an area under a curve computed over this distribution is largely
determined by how ties are broken.}
\label{fig:ecdf}
\end{figure}

Rebuilding the identity over three independently tagged concepts (assets, liabilities
and equity) required one additional extraction pass and produced no usable evaluation,
for a reason that generalises into a procedure. The relation
$\text{assets} = \text{liabilities} + \text{equity}$ holds of a consolidated balance
sheet only once non-controlling and redeemable interests are counted in equity, and the
concept a practitioner reaches for first is parent-company equity, which excludes them.
On our corpus the reference values violate the relation for $47\%$ of filings, always by
the omitted minority interest. Excluding filings whose inputs are not all tagged, and
those whose reference values do not balance, leaves $28$ documents per scale containing
three and one erroneous extraction: too few to estimate anything, with precision
intervals spanning $[0.3, 52.7]$ and $[0, 100]$.

The procedure this suggests costs nothing and would have caught it. Before treating a
relation as a verification signal, evaluate it on the reference values alone. If it does
not hold there, it is not an identity of the standard as one has operationalised it,
whatever its textbook form, and any residual computed from an extraction confounds
misreading with the modelling choice behind the reference. We report the uncoupled
detector as untestable on this corpus rather than refuted, and note that the population
on which it is testable is determined by tagging practice rather than by the detector.

\section{What Survives}
\label{sec:survives}

\begin{figure*}[t]
\centering
\includegraphics[width=0.86\linewidth]{figures/field_sector.pdf}
\caption{Zero-shot accuracy by field and sector (8B, $\pm1\%$). Difficulty is
structured by disclosure convention rather than spread evenly: the banking row is
uniformly weak because commercial-bank income statements carry no operating-income
line, and the blank cell is a concept those filers do not tag at all.}
\label{fig:heat}
\end{figure*}

Zero-shot extraction is accurate and scaffolding does not improve it. Over both
statements and both scales, accuracy at the operating tolerance is $85.9\%$ and
$91.4\%$, and the corrector matches it exactly at every tolerance
(Table~\ref{tab:survives}). Accuracy is tolerance-sensitive: requiring exact agreement
costs a further twelve points on the smaller model, which any deployment threshold
should account for (Figure~\ref{fig:ladder}a).

\begin{table}[t]
\centering\small
\begin{tabular}{@{}lrr@{}}
\toprule
\textbf{Tolerance} & \textbf{8B} & \textbf{49B} \\
\midrule
exact match & 73.9 & 84.6 \\
$\pm0.1\%$  & 76.9 & 86.3 \\
$\pm1\%$    & \best{85.9} & \best{91.4} \\
$\pm5\%$    & 89.5 & 93.9 \\
\bottomrule
\end{tabular}
\caption{Zero-shot accuracy over income-statement and balance-sheet fields, coupled
field excluded, $590$ instances over $81$ documents. The corrector equals these values
at every tolerance.}
\label{tab:survives}
\end{table}

The regex baseline is compared on income-statement fields only, its patterns not
covering the balance sheet. There it reaches $71.7\%$ against $87.0\%$ and $91.0\%$ for
zero-shot extraction, gaps of $15.3$ and $19.3$ points at $p = 6{\cdot}10^{-8}$ and
$2{\cdot}10^{-12}$. That a language model reads these tables better than a pattern
matcher is not itself a finding; it is reported because it is the only comparison in the
study that reaches significance, which calibrates what the null on the remaining
configurations does and does not say. The instrument can detect an effect of this kind
when one is present.

Accuracy varies by sector far more than by configuration (Figure~\ref{fig:heat}).
Banking is weakest on both scales, and it is a schema mismatch rather than a reading
failure: commercial-bank income statements are built from net interest income and
non-interest revenue and carry no operating-income line, so no rule defined over that
quantity can help. Energy fails on revenue and net income for a related reason,
depletion and certain totals being disclosed outside the extracted section. Cells are
small and no per-sector claim is made; the pattern is reported because it is where a
domain-specific schema, rather than a better model or more scaffolding, would help.

Finally, the coupled field is genuinely unenforceable rather than merely outside the
extraction window. Searching the complete text of all $104$ filings, $26.9\%$ mention
EBITDA at all, and $4.8\%$ present it with a reconciliation. The dominant form of the
mention is contractual: EBITDA appears as a defined term in credit-agreement covenants,
not as a disclosed quantity, and no income statement in the corpus reports it as a line
item.

\section{Trap Catalogue}
\label{sec:traps}

Table~\ref{tab:traps} states each trap as a mechanism, the instrument that detects it,
and the evidence observed here. Two of the four are specific to constraint-based
verification; the other two apply to any evaluation that compares extracted quantities
against a reference.

\begin{table*}[t]
\centering\small
\begin{tabular}{@{}p{2.7cm}p{4.3cm}p{3.5cm}p{3.9cm}@{}}
\toprule
\textbf{Trap} & \textbf{Mechanism} & \textbf{Instrument} & \textbf{Evidence here} \\
\midrule
Method-coupled label &
A field's reference value is produced by a formula the system also applies, so accuracy
on it cannot fail &
Rescore with such fields excluded; predict the inflation from the label definition
(\S\ref{sec:predict}) &
$15.5\%$ of scored instances, $100\%$ of the reported gain; reproduced on 15 injected
couplings \\
\addlinespace[2pt]
Ablation by accuracy delta &
Leave-one-out cannot distinguish a redundant rule set from an irrelevant one &
Per-rule execution counts &
Eight rules reported inert; one ever fires \\
\addlinespace[2pt]
Threshold mismatch &
A detector firing on any nonzero residual, judged against a metric that forgives one
percent, manufactures false positives &
Report at matched thresholds, or sweep both (Figure~\ref{fig:surface}) &
Precision $80.0\%$ rises to $100\%$; all four false positives eliminated \\
\addlinespace[2pt]
Unit-scale mismatch &
References in absolute amounts compared against tables printed in millions fail every
comparison while the code runs normally &
Infer scale from the document, never from the reference; report scale as its own error
kind &
$12$ filings excluded for want of a readable scale \\
\bottomrule
\end{tabular}
\caption{The four measurement traps. A scale error is the one failure invisible to this
entire class of method, since scaling every figure by the same factor leaves all
accounting identities exactly satisfied. Deriving a missing label, as in the first row,
is a modelling decision and should be reported as one.}
\label{tab:traps}
\end{table*}

\paragraph{Recommendations.} Verify over quantities the filer is required to disclose,
and confirm the identity holds in the reference data before assuming it holds in the
extraction. Report per-field sample sizes, since tag coverage is a property of the
filer. Report class balance alongside any ranking metric, and prefer a confusion matrix
when the score distribution is dominated by ties.

\section{Conclusion}

An evaluation of accounting-identity verification reported substantial gains that, on
audit, are entirely attributable to one field whose reference value the system
reconstructs by its own formula. Excluding it, six configurations across two model
scales are statistically indistinguishable and four are bit-identical. The detector's
apparent precision has the same origin in a different form, and the uncoupled identity
turns out not to be an identity of the reporting standard. What remains is that
zero-shot extraction is accurate, substantially better than pattern matching, and
improved by none of the scaffolding we tested. We release the corpus manifest, cached
outputs and audit scripts so that these conclusions can themselves be re-evaluated.

\section*{Limitations}
\label{sec:limitations}

\paragraph{Scope of the null.} An exact null over $253$ instances bounds the effect of
these configurations tightly, but the corpus is one language, one filing type and two
model scales from one family on one API, and each technique is instantiated once. Our
chain-of-thought verifier is scoped to a single identity, our self-consistency draws
five samples at one temperature, and our corrector holds eight rules of which one
executes. A differently scoped verifier could plausibly do better, and
\S\ref{sec:collapse} argues that coverage is where the headroom lies. The negative
claim concerns these implementations, not the techniques.

\paragraph{The uncoupled detector is unresolved.} After excluding filings with untagged
inputs and those whose reference data do not balance, $26$ documents remain per scale
with one or two erroneous extractions. This is a statement about evaluability on this
corpus, not evidence that the uncoupled detector fails.

\paragraph{No expert validation.} Every reference value comes from a regulatory tag,
matched by concept name, fiscal period and unit. Neither the tags nor the matching were
checked by a financial professional on any sample, so a systematic mis-mapping would be
invisible to us. The scale inference is validated against observable ratios
(\S\ref{sec:protocol}) and the field-to-concept mapping is released
(Appendix~\ref{app:concepts}), but neither substitutes for expert review, and the
minority-interest case of \S\ref{sec:audit} shows what such a mis-mapping looks like
when it goes unnoticed.

\paragraph{Reference concepts.} Accuracy is accuracy against a specific choice of
concepts. An earlier balance-sheet pass used field names that led one model to
interpret total liabilities as excluding equity; that run is retained as an ablation.
Other fields may carry smaller versions of the same effect.

\paragraph{Extraction window.} Text is sliced around the income statement, so
quantities disclosed only in the cash-flow statement are out of reach, which is a known
cause of the energy-sector result. Filers submitting statements as a separate exhibit
are excluded.

\paragraph{Cell sizes.} Five sectors have four or fewer documents. Per-sector figures
are descriptive; no significance test is reported and none should be inferred.

\paragraph{Not a judgement of the filings.} The system checks arithmetic consistency
and textual grounding. It does not detect misstatement, fraud or accounting-policy
choices.

\section*{Ethics Statement}

All data are public regulatory filings retrieved through official interfaces under the
issuer's fair-access policy with an identifying user agent. No personal, confidential
or proprietary data is used. The specific risk here is over-confidence: a system
reporting a document as verified may be read as warranting its figures, when it checks
only internal arithmetic. This paper's principal finding is that such warrants were
largely circular in our own system; we recommend that identity-based verification be
used for triage rather than sign-off, and that any reported gain be accompanied by the
coupled-field test.

\section*{Acknowledgments}

Removed for anonymous review.

\bibliography{custom}

\appendix

\section{Field-to-Concept Mapping}
\label{app:concepts}

Table~\ref{tab:concepts} gives the mapping. The first concept with a value for the
period wins and the choice is recorded per label, so a reviewer can re-derive any
reference value. The final row is the coupled field: it has no tag, and the derivation
used to fill that gap is what \S\ref{sec:audit} analyses.

\begin{table}[h]
\centering\small
\begin{tabular}{@{}p{1.9cm}p{4.7cm}@{}}
\toprule
\textbf{Field} & \textbf{Concept} \\
\midrule
Revenue & \code{Revenues}, else \code{Revenue\-From\-Contract\-With\-Customer\-Excluding\-Assessed\-Tax} \\
Op.\ income & \code{Operating\-Income\-Loss} \\
D\&A & \code{Depreciation\-Depletion\-And\-Amortization} \\
Net income & \code{Net\-Income\-Loss} \\
Total assets & \code{Assets} \\
Total liab.\ \& eq. & \code{Liabilities\-And\-Stockholders\-Equity} \\
Total liabilities & \code{Liabilities} \\
Total equity & \code{Stockholders\-Equity} \\
\midrule
EBITDA & \emph{no tag exists} --- derived as
         op.\ income ${+}$ D\&A \\
\bottomrule
\end{tabular}
\caption{Field-to-concept mapping.}
\label{tab:concepts}
\end{table}

\section{Corpus Reconciliation}
\label{app:corpus}

\begin{table}[h]
\centering\small
\begin{tabular}{@{}lr@{}}
\toprule
\textbf{Step} & \textbf{Documents} \\
\midrule
62 tickers $\times$ 2 filings & 124 \\
$-$ filings absent from the recent window & $-8$ \\
$-$ no readable reporting scale & $-12$ \\
\midrule
\textbf{corpus} & \textbf{104} \\
\bottomrule
\end{tabular}
\caption{The eight absent filings are six banks with a single recent annual report and
one company with none. The twelve scale exclusions are six companies whose statements
are filed as a separate exhibit, both of their years. Per-document reasons are released
with the corpus.}
\label{tab:corpusrec}
\end{table}

Ground-truth coverage differs by field, which is why per-field $n$ is reported
throughout: net income, total assets, total liabilities and equity, and total equity
are tagged for every document; revenue for $97\%$; operating income and D\&A for
$81\%$; total liabilities for $65\%$; and EBITDA for none.

\section{Populations Used by Each Table}
\label{app:populations}

Results are reported on several nested subsets, because coverage differs by extraction
pass, by tag availability and by model. Table~\ref{tab:populations} states which subset
each result uses and why it differs from the full corpus. No subset is chosen by
outcome: each is the set of documents for which the required inputs exist.

\begin{table}[h]
\centering\small
\begin{tabular}{@{}p{2.1cm}rp{4.2cm}@{}}
\toprule
\textbf{Result} & \textbf{$n$} & \textbf{Population and reason} \\
\midrule
Corpus (\S\ref{sec:protocol}) & 104 &
All filings with a readable reporting scale and matching reference values. \\
\addlinespace[2pt]
Table~\ref{tab:asfirst}, \ref{tab:rescore} & 83 &
Filings both model scales completed for the income-statement pass. \\
\addlinespace[2pt]
Table~\ref{tab:mcnemar} & 83 &
The same set; B4 is compared on the subset it returned. \\
\addlinespace[2pt]
Table~\ref{tab:fires} & 83 &
Corrector executions, income-statement pass only. \\
\addlinespace[2pt]
Table~\ref{tab:survives} & 81 &
Filings both scales completed for the income-statement \emph{and} the four-field
balance-sheet pass. \\
\addlinespace[2pt]
Fig.~\ref{fig:heat}, per-sector & 81 &
As above, over the three balance-sheet fields common to both balance-sheet passes. \\
\addlinespace[2pt]
Coupled detector (\S\ref{sec:audit}) & 68 &
Filings with a computable residual on the two-field balance-sheet pass. \\
\addlinespace[2pt]
Uncoupled detector & 28 &
Of those, filings whose three identity inputs are all tagged and whose reference values
satisfy the identity. \\
\bottomrule
\end{tabular}
\caption{Population per result. The two detector rows are the only large reductions, and
both follow from tagging practice: total liabilities is tagged by $65\%$ of companies,
and the reference values satisfy
$\text{assets}=\text{liabilities}+\text{equity}$ for $53\%$ of the remainder.}
\label{tab:populations}
\end{table}

Rate limiting on the larger model prevented completing every pass on all 104 filings.
A third model was extracted for 47 filings and discarded before analysis: its subset
contained 28 technology filings and no retail, consumer or industrial ones, so a tier
comparison on it would have confounded model scale with sector. That decision was taken
on the composition of the subset, not on any result computed from it.

\section{Detector Confusion Matrices}
\label{app:detector}

\begin{table}[h]
\centering\small
\begin{tabular}{@{}lrrrr@{}}
\toprule
\textbf{Threshold} & \textbf{TP} & \textbf{FP} & \textbf{FN} & \textbf{P (95\% CI)} \\
\midrule
residual $>0$      & 16 & 4 & 3 & 80.0 $[56.3, 94.3]$ \\
residual $>0.1\%$  & 16 & 3 & 3 & 84.2 $[60.4, 96.6]$ \\
residual $>1\%$    & 16 & 0 & 3 & \best{100.0} $[79.4, 100]$ \\
\bottomrule
\end{tabular}
\caption{Coupled balance-sheet detector, 8B, $68$ documents, positive label defined as
at least one identity input differing from its reference by more than $1\%$. Matching
the detector threshold to the correctness tolerance eliminates every false positive.}
\label{tab:detconf}
\end{table}

The four false positives at the loosest threshold are documents whose largest field
error is between $0.004\%$ and $0.35\%$, all inside the correctness tolerance.
Figure~\ref{fig:surface} generalises this to the full threshold plane.

\section{Extraction Prompt}
\label{app:prompts}

Identical across B2--B6 and both scales; only the scaffolding differs.

\begin{quote}\footnotesize\ttfamily\raggedright
You are reading the consolidated income statement of a US 10-K filing. Extract these 5
metrics for the MOST RECENT fiscal year shown.\\[2pt]
Return ONLY a JSON object with exactly these keys (use null when absent):
revenue, operating\_income, ebitda, depreciation\_\-amortization, net\_income\\[2pt]
Rules: Copy each number EXACTLY as printed in the table. Do NOT rescale or convert
units. Consolidated group totals only, never a single segment. Most recent fiscal year
column only, never a comparative year. No markdown, no commentary. JSON only.
\end{quote}

Two instructions in this prompt interact with the audit. The instruction not to rescale
is what makes document-side scale inference necessary. The instruction to return null
when a value is absent is what the coupled field then penalised, since EBITDA is absent
from every income statement in the corpus and the model complied throughout.

\section{Reproducibility}
\label{app:repro}

All tables and figures are recomputable offline from the released artifacts, with no
model access.

\begin{itemize}[leftmargin=*, itemsep=1pt]
\item \code{data/corpus.jsonl} --- manifest with path, ticker, identifier, fiscal year,
      period end, accession number and inferred scale; checksummed. Drop reasons for
      every excluded document are released alongside.
\item \code{data/ground\_truth.jsonl} --- reference values with the concept, accession
      number and period behind each, plus a \code{derived} flag marking coupled labels.
\item \code{data/extractions/<model>/<method>/} --- every cached model output.
\item \code{scripts/reviewer\_checks.py} --- coupled-field rescoring and per-rule
      execution counts.
\item \code{scripts/coupling\_index.py} --- the predictor of \S\ref{sec:predict}.
\item \code{scripts/mcnemar\_ladder.py} --- discordant pairs and exact tests.
\item \code{scripts/detector\_and\_bs.py} --- confusion matrices, matched thresholds,
      tolerance sweep.
\item \code{scripts/make\_figures.py} --- all figures.
\end{itemize}

Seed fixed at $42$; temperature $0$ except B3 ($T{=}0.7$, five samples); $60$k
characters of context. Failed API calls are never cached as null results, since a
cached null is indistinguishable from a model that extracted nothing. An early run
cached $48$ of $104$ documents that way, which would have depressed all configurations
equally and resembled a model result.

\end{document}
